% TOP_PROBLEM: {"id": 3151, "title": "Graham's conjecture on tree reconstruction", "category_id": 3, "category": {"id": 3, "name": "graph_theory", "display_name": "Graph Theory", "description": "Problems involving graphs, networks, and their properties.", "slug": "graph-theory", "order_index": 3, "created_at": "2026-07-31T15:26:25.671Z"}, "status": "open", "problem_number": "OPG-164", "set_id": 12}
% FIRST_POSTED: 2026-10-01
% ENTRY_KIND: statement_only
% UnsolvedMath ID 3151; source catalogue code OPG-164
% !TeX program = lualatex
% Definition and sources only; this file is not an LLM solution attempt.
\documentclass[11pt,a4paper]{article}
\usepackage[margin=25mm]{geometry}
\usepackage{fontspec}
\setmainfont{lmroman10-regular.otf}[BoldFont=lmroman10-bold.otf,
 ItalicFont=lmroman10-italic.otf,BoldItalicFont=lmroman10-bolditalic.otf]
\usepackage{amsmath,amssymb,mathtools}
\usepackage{unicode-math}
\setmathfont{latinmodern-math.otf}
\AtBeginDocument{\let\setminus\smallsetminus}
\usepackage{xurl,hyperref}
\hypersetup{colorlinks=true,urlcolor=blue}
\setcounter{secnumdepth}{0}
\setlength{\parindent}{0pt}
\setlength{\parskip}{5pt}
\setlength{\emergencystretch}{3em}
\begin{document}
\section*{Graham's conjecture on tree reconstruction}\label{3151}
{\small UnsolvedMath ID 3151 · OPG-164 · Graph Theory · OpenGarden}
\subsection{Definitions and mathematical statement}
For a finite simple graph $G$, its line graph $L(G)$ has one vertex
for each edge of $G$. Two distinct vertices of $L(G)$ are adjacent
precisely when the corresponding edges of $G$ share an endpoint.
Define iterates by $L^0(G)=G$ and $L^{j+1}(G)=L(L^j(G))$ for every
integer $j\geq0$. Each iterate is again a finite simple graph.
The line graph of an edgeless graph has no vertices, and subsequent
iterates stay empty.

Associate with $G$ its infinite sequence of orders
\[
 a(G)=(a_0(G),a_1(G),\ldots),\qquad a_j(G)=|V(L^j(G))|.
\]
A tree is a nonempty connected graph with no cycles. Graham's
reconstruction question is whether
\[
 \bigl[a_j(T)=a_j(T')\text{ for every }j\geq0\bigr]
 \quad\Longrightarrow\quad T\cong T'
\]
for every pair of finite trees. Isomorphism allows any renaming
of vertices and must preserve the edge relation.

Only the numbers of vertices are supplied, not the line graphs
themselves, their degree sequences, or maps between successive
iterates. The first two entries of an $n$-vertex tree are $n$ and
$n-1$, while later entries reflect more structure. The entire
infinite sequence is part of the statement: a collision among
finitely many initial entries alone does not disprove the
conjecture. A negative answer requires nonisomorphic trees with
agreement at every iteration.
\subsection{Short English statement}
Can a finite tree be reconstructed, up to isomorphism, from the vertex counts of all its iterated line graphs?
\subsection{Sources}
\begin{itemize}
\item[S1] ULAM.AI and the UnsolvedMath Contributors, supplied problems.json, record 3151 (OPG-164), OpenGarden collection. Catalogue identity and source transcription; CC BY 4.0.
\url{https://huggingface.co/datasets/ulamai/UnsolvedMath}.
\item[S2] Open Problem Garden, Graham's conjecture on tree reconstruction. Original problem and discussion; the mathematical formulation below is expanded from the supplied source record.
\url{http://www.openproblemgarden.org/op/grahams_conjecture_on_tree_reconstruction}.
\end{itemize}
\end{document}
